\documentclass[11pt,a4paper]{article}

\usepackage[utf8]{inputenc}
\usepackage[T1]{fontenc}
\usepackage[spanish]{babel}
\usepackage{amsmath,amssymb,mathtools}
\usepackage{geometry}
\usepackage{lmodern}
\usepackage{microtype}

\geometry{margin=2.4cm}
\setlength{\parindent}{0pt}
\setlength{\parskip}{0.45em}

\newcommand{\R}{\mathbb{R}}

\begin{document}

\begin{center}
{\LARGE\bfseries Geometría Lineal}\\[0.4em]
{\Large Tema 1: Fundamentos de geometría plana}\\[0.4em]
{\large Soluciones}
\end{center}

\section{Ejercicio 1}
\subsection{Solución}

Los lados del triángulo son:

\[
15,\qquad 20,\qquad 25.
\]

Como:

\[
15^2+20^2
=
225+400
=
625
=
25^2,
\]

el triángulo es rectángulo.

Por tanto:

\[
A
=
\frac12\cdot15\cdot20.
\]

Así:

\[
\boxed{A=150}
\]

\section{Ejercicio 2}
\subsection{Solución}

El cuadrado tiene lado:

\[
l=\sqrt{100}=10.
\]

Sea \(x\) el cateto de cada triángulo rectángulo isósceles eliminado.

Los lados inclinados del octógono miden:

\[
x\sqrt2.
\]

Los lados horizontales y verticales miden:

\[
10-2x.
\]

Como el octógono es regular:

\[
x\sqrt2=10-2x.
\]

Por tanto:

\[
x
=
\frac{10}{2+\sqrt2}.
\]

El lado del octógono es:

\[
s=x\sqrt2
=
10(\sqrt2-1).
\]

Así, su perímetro es:

\[
\boxed{
P
=
80(\sqrt2-1)
}
\]

\section{Ejercicio 3}
\subsection{Solución}

La circunferencia tiene radio:

\[
R=2.
\]

Por tanto:

\[
A_C=4\pi.
\]

El hexágono regular tiene lado:

\[
l=1.
\]

Su área es:

\[
A_H
=
6\frac{\sqrt3}{4}
=
\frac{3\sqrt3}{2}.
\]

La región morada es el área de la circunferencia menos la del hexágono:

\[
\boxed{
A
=
4\pi-\frac{3\sqrt3}{2}
}
\]

\section{Ejercicio 4}
\subsection{Solución}

Sea:

\[
AB=CD=x.
\]

Como:

\[
BC=12,
\qquad
BE=CE=10,
\]

la altura desde \(E\) al punto medio de \(BC\) es:

\[
h
=
\sqrt{10^2-6^2}
=
8.
\]

Además:

\[
AD=2x+12.
\]

Por simetría:

\[
AE=ED
=
\sqrt{(x+6)^2+8^2}.
\]

El perímetro de \(\triangle BEC\) es:

\[
10+10+12=32.
\]

Por tanto:

\[
2\sqrt{(x+6)^2+64}
+
2x+12
=
64.
\]

Dividiendo entre \(2\):

\[
\sqrt{(x+6)^2+64}
+
x+6
=
32.
\]

Sea:

\[
t=x+6.
\]

Entonces:

\[
\sqrt{t^2+64}=32-t.
\]

Elevando al cuadrado:

\[
t^2+64
=
1024-64t+t^2.
\]

Así:

\[
t=15.
\]

Por tanto:

\[
\boxed{AB=9}
\]

\section{Ejercicio 5}
\subsection{Solución}

El lado del cuadrado exterior es:

\[
AB=\sqrt{50}.
\]

Por tanto, su diagonal mide:

\[
DB
=
\sqrt{50}\sqrt2
=
10.
\]

El centro del cuadrado divide la diagonal en dos segmentos de longitud:

\[
5.
\]

Como:

\[
BE=1,
\]

la distancia desde \(E\) al centro es:

\[
5-1=4.
\]

El cuadrado \(EFGH\) tiene por diagonal:

\[
EG=8.
\]

El área de un cuadrado de diagonal \(d\) es:

\[
A=\frac{d^2}{2}.
\]

Por tanto:

\[
\boxed{
A_{EFGH}
=
\frac{8^2}{2}
=
32
}
\]

\section{Ejercicio 6}
\subsection{Solución}

Sea \(h\) la altura y sean \(x,y\) las proyecciones horizontales de los lados de longitudes \(5\) y \(12\).

Como las bases miden \(52\) y \(39\):

\[
x+y=13.
\]

Además:

\[
x^2+h^2=25,
\]

\[
y^2+h^2=144.
\]

Restando:

\[
y^2-x^2=119.
\]

Así:

\[
(y-x)(y+x)=119.
\]

Como:

\[
x+y=13,
\]

se obtiene:

\[
y-x=\frac{119}{13}.
\]

Por tanto:

\[
x=\frac{25}{13},
\qquad
y=\frac{144}{13}.
\]

Entonces:

\[
h
=
\sqrt{
25-\left(\frac{25}{13}\right)^2
}
=
\frac{60}{13}.
\]

El área del trapecio es:

\[
A
=
\frac{52+39}{2}
\frac{60}{13}.
\]

Por tanto:

\[
\boxed{A=210}
\]

\section{Ejercicio 7}
\subsection{Solución}

El ángulo central de un octógono regular es:

\[
\frac{360^\circ}{8}=45^\circ.
\]

Si el lado mide \(2\), su radio circunscrito es:

\[
R
=
\frac{2}{2\sen(22.5^\circ)}.
\]

Tomando los vértices \(A,D,G\) sobre la circunferencia, los ángulos centrales correspondientes son:

\[
\angle AOD=135^\circ,
\]

\[
\angle DOG=135^\circ,
\]

\[
\angle GOA=90^\circ.
\]

Aplicando la fórmula del área mediante coordenadas, o equivalentemente sumando áreas de triángulos con vértice en el centro, se obtiene:

\[
A_{\triangle ADG}
=
\frac{R^2}{2}
\left(
\sen135^\circ
+
\sen135^\circ
+
\sen90^\circ
\right).
\]

Como:

\[
R^2=4+2\sqrt2,
\]

resulta:

\[
\boxed{
A_{\triangle ADG}
=
4+3\sqrt2
}
\]

\section{Ejercicio 8}
\subsection{Solución}

Sean:

\[
OX=a,
\qquad
OY=b.
\]

Como \(M\) y \(N\) son puntos medios:

\[
OM=\frac a2,
\qquad
ON=\frac b2.
\]

Por Pitágoras:

\[
XN^2
=
a^2+\frac{b^2}{4}
=
19^2,
\]

\[
YM^2
=
b^2+\frac{a^2}{4}
=
22^2.
\]

Multiplicando por \(4\):

\[
4a^2+b^2=1444,
\]

\[
a^2+4b^2=1936.
\]

Resolviendo:

\[
a^2=256,
\qquad
b^2=420.
\]

Por tanto:

\[
XY^2
=
a^2+b^2
=
676.
\]

Así:

\[
\boxed{XY=26}
\]

\section{Ejercicio 9}
\subsection{Solución}

El semicírculo pequeño tiene diámetro \(1\), por lo que su radio es:

\[
r=\frac12.
\]

Su área es:

\[
A_1
=
\frac12\pi\left(\frac12\right)^2
=
\frac{\pi}{8}.
\]

El semicírculo grande tiene radio:

\[
R=1.
\]

La cuerda común mide \(1\), luego el ángulo central correspondiente es:

\[
60^\circ=\frac{\pi}{3}.
\]

El segmento circular situado sobre la cuerda tiene área:

\[
A_2
=
\frac{\pi}{6}
-
\frac{\sqrt3}{4}.
\]

Por tanto, la región morada mide:

\[
A
=
A_1-A_2.
\]

Así:

\[
\boxed{
A
=
\frac{\sqrt3}{4}
-
\frac{\pi}{24}
}
\]

\section{Ejercicio 10}
\subsection{Solución}

Tomamos:

\[
A=(0,0),
\qquad
B=(4,0),
\]

\[
E=(0,8),
\qquad
C=(4,6).
\]

Las rectas \(AC\) y \(BE\) son:

\[
AC:
\quad
y=\frac32x,
\]

\[
BE:
\quad
y=-2x+8.
\]

Su intersección es:

\[
D=
\left(
\frac{16}{7},
\frac{24}{7}
\right).
\]

El área de \(\triangle ADE\) es:

\[
A_{ADE}
=
\frac12\cdot8\cdot\frac{16}{7}
=
\frac{64}{7}.
\]

El área de \(\triangle BDC\) es:

\[
A_{BDC}
=
\frac12\cdot6
\left(
4-\frac{16}{7}
\right)
=
\frac{36}{7}.
\]

Por tanto:

\[
\boxed{
A_{ADE}-A_{BDC}=4
}
\]

\section{Ejercicio 11}
\subsection{Solución}

El semicírculo grande tiene radio \(2\), luego:

\[
A_G
=
\frac12\pi(2)^2
=
2\pi.
\]

Cada semicírculo pequeño tiene radio \(1\), por lo que su área es:

\[
\frac{\pi}{2}.
\]

Los centros consecutivos están separados una unidad.

La intersección de dos circunferencias unitarias cuyos centros distan \(1\) tiene área:

\[
2\frac{\pi}{3}
-
\frac{\sqrt3}{2}.
\]

La mitad superior de dicha intersección mide:

\[
\frac{\pi}{3}
-
\frac{\sqrt3}{4}.
\]

Hay dos solapamientos de este tipo, por lo que el área ocupada por los tres semicírculos pequeños es:

\[
A_P
=
3\frac{\pi}{2}
-
2\left(
\frac{\pi}{3}
-
\frac{\sqrt3}{4}
\right).
\]

Así:

\[
A_P
=
\frac{5\pi}{6}
+
\frac{\sqrt3}{2}.
\]

Por tanto:

\[
\boxed{
A
=
2\pi-A_P
=
\frac{7\pi}{6}
-
\frac{\sqrt3}{2}
}
\]

\section{Ejercicio 12}
\subsection{Solución}

Tomamos:

\[
AB=BC=CD=DE=EF=1.
\]

Entonces:

\[
AD=3,
\qquad
AF=5.
\]

Como:

\[
HC\parallel AG,
\]

los triángulos:

\[
\triangle DCH
\quad\text{y}\quad
\triangle DAG
\]

son semejantes.

Por tanto:

\[
\frac{HC}{AG}
=
\frac{DC}{DA}
=
\frac13.
\]

Además:

\[
JE\parallel AG,
\]

por lo que:

\[
\triangle FEJ
\sim
\triangle FAG.
\]

Así:

\[
\frac{JE}{AG}
=
\frac{FE}{FA}
=
\frac15.
\]

Por tanto:

\[
\boxed{
\frac{HC}{JE}
=
\frac{1/3}{1/5}
=
\frac53
}
\]

\section{Ejercicio 13}
\subsection{Solución}

Tomamos:

\[
D=(0,0),
\qquad
A=(8,0),
\]

\[
C=(0,9),
\qquad
B=(8,9).
\]

Como:

\[
BH=6,
\]

se tiene:

\[
H=(2,9).
\]

Como:

\[
DE=4,
\]

se tiene:

\[
E=(4,0).
\]

La recta \(AH\) es:

\[
y=-\frac32(x-8),
\]

y la recta \(EC\):

\[
y=-\frac94(x-4).
\]

Resolviendo ambas ecuaciones:

\[
G=(-4,18).
\]

Como:

\[
GF\perp AF,
\]

la recta \(GF\) es vertical y \(F=(-4,0)\).

Por tanto:

\[
\boxed{GF=18}
\]

\section{Ejercicio 14}
\subsection{Solución}

Tomamos:

\[
D=(0,0),
\qquad
C=(5,0),
\]

\[
A=(0,3),
\qquad
B=(5,3).
\]

Como:

\[
DF=1,
\]

\[
F=(1,0).
\]

Como:

\[
GC=2,
\]

\[
G=(3,0).
\]

Las rectas \(AF\) y \(BG\) se cortan en:

\[
E=
\left(
\frac53,-2
\right).
\]

El segmento \(AB\) tiene longitud \(5\) y la distancia de \(E\) a la recta \(AB\) es:

\[
5.
\]

Por tanto:

\[
A_{\triangle AEB}
=
\frac12\cdot5\cdot5.
\]

Así:

\[
\boxed{
A_{\triangle AEB}
=
\frac{25}{2}
}
\]

\section{Ejercicio 15}
\subsection{Solución}

Tomamos un cuadrado de lado \(1\):

\[
D=(0,0),
\qquad
C=(1,0),
\]

\[
A=(0,1),
\qquad
B=(1,1).
\]

Sean:

\[
E=(0,t),
\qquad
F=(t,0).
\]

Como \(\triangle BEF\) es equilátero:

\[
BE=BF=EF.
\]

De:

\[
BE^2=EF^2,
\]

se obtiene:

\[
1+(1-t)^2=2t^2.
\]

Por tanto:

\[
t=\sqrt3-1.
\]

El área de \(\triangle DEF\) es:

\[
A_{DEF}
=
\frac{t^2}{2}.
\]

El área de \(\triangle ABE\) es:

\[
A_{ABE}
=
\frac{1-t}{2}.
\]

Así:

\[
\frac{A_{DEF}}{A_{ABE}}
=
\frac{t^2}{1-t}.
\]

Sustituyendo:

\[
t=\sqrt3-1,
\]

se obtiene:

\[
\boxed{
\frac{A_{DEF}}{A_{ABE}}=2
}
\]

\section{Ejercicio 16}
\subsection{Solución}

Sea \(\alpha\) el ángulo agudo que forman las dos rectas.

En el círculo de radio \(1\), la región sombreada corresponde a los dos sectores agudos, por lo que su área es:

\[
A_1=\alpha.
\]

En la corona comprendida entre radios \(1\) y \(2\), la región sombreada corresponde a los sectores obtusos:

\[
A_2
=
3(\pi-\alpha).
\]

En la corona comprendida entre radios \(2\) y \(3\), la región sombreada corresponde de nuevo a los sectores agudos:

\[
A_3
=
5\alpha.
\]

Así:

\[
A_S
=
\alpha+3(\pi-\alpha)+5\alpha
=
3\pi+3\alpha.
\]

El área total es:

\[
9\pi.
\]

Por tanto, el área no sombreada es:

\[
A_N
=
6\pi-3\alpha.
\]

Como:

\[
A_S
=
\frac8{13}A_N,
\]

se tiene:

\[
3\pi+3\alpha
=
\frac8{13}
(6\pi-3\alpha).
\]

Resolviendo:

\[
\alpha=\frac{\pi}{7}.
\]

Por tanto:

\[
\boxed{
\alpha
=
\frac{\pi}{7}
=
\frac{180^\circ}{7}
}
\]

\section{Ejercicio 17}
\subsection{Solución}

Sea \(O\) el centro de la circunferencia \(D\).

La circunferencia \(A\) tiene radio \(1\) y pasa por \(O\), por lo que la distancia entre su centro y \(O\) es:

\[
1.
\]

Como además \(A\) es tangente interiormente a \(D\):

\[
R_D-1=1.
\]

Así:

\[
R_D=2.
\]

Sea \(r\) el radio común de \(B\) y \(C\).

Por simetría, tomamos sus centros como:

\[
(x,r),
\qquad
(x,-r).
\]

Como son tangentes interiormente a \(D\):

\[
x^2+r^2=(2-r)^2.
\]

Por tanto:

\[
x^2=4-4r.
\]

Como son tangentes exteriormente a \(A\), cuyo centro podemos tomar en \((1,0)\):

\[
(x-1)^2+r^2=(1+r)^2.
\]

Simplificando:

\[
x=2-3r.
\]

Sustituyendo:

\[
(2-3r)^2=4-4r.
\]

Así:

\[
9r^2-8r=0.
\]

Como \(r>0\):

\[
\boxed{
r=\frac89
}
\]

\section{Ejercicio 18}
\subsection{Solución}

Tomamos el triángulo rectángulo:

\[
C=(0,0),
\qquad
A=(12,0),
\qquad
E=(0,16).
\]

Como:

\[
AB=3,
\]

se tiene:

\[
B=(9,0).
\]

Como:

\[
CD=4,
\]

\[
D=(0,4).
\]

Además:

\[
EA=20,
\qquad
EF=5.
\]

Por tanto, \(F\) divide \(EA\) en razón:

\[
\frac{EF}{EA}=\frac14.
\]

Así:

\[
F=(3,12).
\]

El área de \(\triangle ACE\) es:

\[
A_{ACE}
=
\frac12\cdot12\cdot16
=
96.
\]

El área de \(\triangle DBF\) es:

\[
A_{DBF}=42.
\]

Por tanto:

\[
\boxed{
\frac{A_{DBF}}{A_{ACE}}
=
\frac{42}{96}
=
\frac7{16}
}
\]

\section{Ejercicio 19}
\subsection{Solución}

Tomamos la circunferencia pequeña de radio \(1\) centrada en el origen.

Sus puntos de tangencia son:

\[
A=(-1,0),
\qquad
B=(1,0).
\]

Las dos circunferencias grandes tienen radio \(2\), y sus centros son:

\[
(1,0),
\qquad
(-1,0).
\]

La distancia entre sus centros es:

\[
d=2.
\]

El área común a dos circunferencias de radio \(R=2\) cuyos centros distan \(2\) es:

\[
A_I
=
2R^2
\arccos\left(
\frac{d}{2R}
\right)
-
\frac d2
\sqrt{4R^2-d^2}.
\]

Por tanto:

\[
A_I
=
8\arccos\left(\frac12\right)
-
\sqrt{12}.
\]

Así:

\[
A_I
=
\frac{8\pi}{3}
-
2\sqrt3.
\]

La circunferencia pequeña está contenida en ambas y tiene área:

\[
\pi.
\]

Por tanto:

\[
\boxed{
A
=
\frac{5\pi}{3}
-
2\sqrt3
}
\]

\section{Ejercicio 20}
\subsection{Solución}

Sean las medianas \(AD\), \(BE\) y \(CF\), que se cortan en el baricentro \(G\).

Como \(D\) es el punto medio de \(BC\):

\[
BD=DC.
\]

Los triángulos:

\[
\triangle BGD
\quad\text{y}\quad
\triangle CGD
\]

tienen la misma altura respecto de la recta \(BC\) y bases iguales.

Por tanto:

\[
A_{BGD}=A_{CGD}.
\]

De forma análoga:

\[
A_{CGE}=A_{AGE}
\]

y:

\[
A_{AGF}=A_{BGF}.
\]

Además, cada mediana divide al triángulo original en dos regiones de igual área.

Combinando estas igualdades se obtiene que los seis triángulos tienen la misma área:

\[
\boxed{
A_{AGF}
=
A_{BGF}
=
A_{BGD}
=
A_{CGD}
=
A_{CGE}
=
A_{AGE}
=
\frac16A_{ABC}
}
\]

\section{Ejercicio 21}
\subsection{Solución}

Sea \(ABCD\) un cuadrilátero inscrito en una circunferencia.

El ángulo:

\[
\angle A
\]

subtiende el arco \(BCD\), por lo que:

\[
\angle A
=
\frac12
m(\widehat{BCD}).
\]

El ángulo opuesto:

\[
\angle C
\]

subtiende el arco \(DAB\), luego:

\[
\angle C
=
\frac12
m(\widehat{DAB}).
\]

Los dos arcos completan la circunferencia:

\[
m(\widehat{BCD})
+
m(\widehat{DAB})
=
360^\circ.
\]

Por tanto:

\[
\angle A+\angle C
=
\frac12\cdot360^\circ.
\]

Así:

\[
\boxed{
\angle A+\angle C=180^\circ
}
\]

Análogamente:

\[
\boxed{
\angle B+\angle D=180^\circ
}
\]

\section{Ejercicio 22}
\subsection{Solución}

Sea \(P\) el punto de intersección de las cuerdas \(AC\) y \(DB\).

Según la figura:

\[
\angle CPB=40^\circ.
\]

Por el teorema de las cuerdas que se cortan en el interior:

\[
40^\circ
=
\frac12
\left(
m(\widehat{CB})
+
m(\widehat{AD})
\right).
\]

Por tanto:

\[
m(\widehat{CB})
+
m(\widehat{AD})
=
80^\circ.
\]

Como \(AB\) es un diámetro, el arco superior \(ADBCA\) mide:

\[
180^\circ.
\]

Así:

\[
m(\widehat{DC})
=
180^\circ-80^\circ
=
100^\circ.
\]

Los ángulos \(\alpha\) y \(\beta\) son ángulos inscritos que subtienen el arco \(DC\).

Por tanto:

\[
\alpha=\beta=50^\circ.
\]

Así:

\[
\boxed{
\alpha+\beta=100^\circ
}
\]

\section{Ejercicio 23}
\subsection{Solución}

En la figura:

\[
AC
\]

y:

\[
BD
\]

son diámetros.

El ángulo:

\[
\angle ACD=60^\circ
\]

subtiende el arco \(AD\).

Por tanto:

\[
m(\widehat{AD})=120^\circ.
\]

Como \(AC\) es un diámetro:

\[
m(\widehat{ADC})=180^\circ.
\]

Así:

\[
m(\widehat{DC})
=
180^\circ-120^\circ
=
60^\circ.
\]

El ángulo \(\alpha=\angle DBC\) subtiende el arco \(DC\), luego:

\[
\boxed{
\alpha
=
\frac12\cdot60^\circ
=
30^\circ
}
\]

\section{Ejercicio 24}
\subsection{Solución}

En la figura:

\[
AB
\]

es un diámetro y:

\[
\angle AOC=140^\circ.
\]

Por tanto:

\[
m(\widehat{AC})=140^\circ.
\]

Como el semicírculo \(ACB\) mide:

\[
180^\circ,
\]

se tiene:

\[
m(\widehat{CB})
=
180^\circ-140^\circ
=
40^\circ.
\]

El ángulo:

\[
\alpha=\angle CDB
\]

es un ángulo inscrito que subtiende el arco \(CB\).

Por tanto:

\[
\boxed{
\alpha
=
\frac12\cdot40^\circ
=
20^\circ
}
\]

\end{document}
